% Source: https://artofproblemsolving.com/community/c3h2455850_help_with_a_cube_question
% Stable identity: aops-topic-2455850-q01
Each column has $5$ cubes and the three columns share only the centre cube, so $5+5+5-2=13$ cubes are removed.

\begin{asy}
usepackage("amsmath");
size(4.2cm);
defaultpen(fontsize(10pt));
pair pr(real x,real y,real z){return (x+0.45*y,z+0.32*y);}
pen hole=rgb(0.35,0.35,0.35);
fill(pr(2,0,2)--pr(3,0,2)--pr(3,0,3)--pr(2,0,3)--cycle,hole);
fill(pr(2,2,5)--pr(3,2,5)--pr(3,3,5)--pr(2,3,5)--cycle,hole);
fill(pr(5,2,2)--pr(5,3,2)--pr(5,3,3)--pr(5,2,3)--cycle,hole);
for(int i=0;i<=5;++i){
 draw(pr(i,0,0)--pr(i,0,5),gray); draw(pr(0,0,i)--pr(5,0,i),gray);
 draw(pr(i,0,5)--pr(i,5,5),gray); draw(pr(0,i,5)--pr(5,i,5),gray);
 draw(pr(5,i,0)--pr(5,i,5),gray); draw(pr(5,0,i)--pr(5,5,i),gray);}
draw(pr(0,0,0)--pr(5,0,0)--pr(5,0,5)--pr(0,0,5)--cycle,linewidth(0.9));
draw(pr(0,0,5)--pr(0,5,5)--pr(5,5,5)--pr(5,0,5),linewidth(0.9));
draw(pr(5,0,0)--pr(5,5,0)--pr(5,5,5),linewidth(0.9));
\end{asy}

\begin{asy}
usepackage("amsmath");
size(3cm);
defaultpen(fontsize(10pt));
pen hole=rgb(0.35,0.35,0.35);
for(int i=0;i<5;++i) for(int j=0;j<5;++j){
  path q=box((i,j),(i+1,j+1));
  if(i==2 || j==2) fill(q,hole);
  draw(q,gray);}
draw(box((0,0),(5,5)),linewidth(0.9));
label("middle layer",(2.5,0),S,fontsize(9pt));
\end{asy}

\textbf{Outer surface.} Each face of the large cube loses the central unit square (where a column comes out), leaving $25-1=24$ unit squares. That gives $6\cdot24=144$.

\textbf{Inner surface.} Each column consists of the centre cube and two ``arms'' of $2$ cubes each, joining the centre cube to two opposite faces; altogether there are $6$ arms of $2$ cubes. In the middle layer (second sketch) four of the arms form a plus sign. The walls of an arm are its $4$ sides, each made of $2$ unit squares, so $8$ unit squares per arm. The ends of each arm are open: one end is a hole in the outer face, and the other end opens into the centre cube. The centre cube has all $6$ of its faces open to arms, so it adds no walls. The inner surface is therefore $6\cdot8=48$.

The total surface area is $144+48=\boxed{192}$.
